\section{Player-to-Game Prediction via Feature Engineering}
\label{sec:player-to-game}

While the state-space model provides robust game-level predictions using team ratings, we hypothesize that incorporating player-level information can improve forecast accuracy. This section explores whether aggregating player statistics into team-level features enhances game outcome prediction---a direct response to the limitations uncovered in the GNN approach.

\subsection{Motivation}

\subsubsection{Why Player-Level Data Matters}

NFL games are fundamentally determined by individual player performance across multiple dimensions:

\begin{enumerate}
    \item \textbf{Quarterback performance}: Passing yards, completion percentage, touchdown rate
    \item \textbf{Running game}: Rushing yards, yards per carry, explosive plays
    \item \textbf{Receiving corps}: Receiving yards, targets, catch rate
    \item \textbf{Star power}: Concentration of elite talent vs.\ depth quality
    \item \textbf{Matchup-specific advantages}: QB vs.\ opponent secondary, RB vs.\ run defense
\end{enumerate}

The state-space model treats teams as monolithic entities with scalar ratings, ignoring this positional heterogeneity. However, as demonstrated in Section~\ref{sec:gnn-failure}, directly modeling player interactions via GNNs leads to severe overfitting (45\% test accuracy, worse than random).

\subsubsection{The Feature Engineering Compromise}

Rather than complex graph structures, we adopt a simpler approach:

\begin{quote}
\textit{Aggregate player-level statistics into interpretable team-level features, then train a lightweight classifier.}
\end{quote}

This balances:
\begin{itemize}
    \item \textbf{Signal}: Leverages rich player data
    \item \textbf{Simplicity}: Avoids overparameterization
    \item \textbf{Interpretability}: Transparent feature importance
\end{itemize}

\subsection{Data Sources}

\subsubsection{NFLverse Player Statistics}

We use the \texttt{nfl\_data\_py} package to access comprehensive player statistics from nflverse \citep{nflverse2023}:

\begin{itemize}
    \item \textbf{Weekly player data}: Per-game statistics for all offensive skill positions (QB, RB, WR, TE)
    \item \textbf{Seasons}: 2019--2024 (2025 data not yet available at time of analysis)
    \item \textbf{Coverage}: 33,287 player-game records across 1,381 unique players
\end{itemize}

\subsubsection{Bayesian Player Projections}

Where available, we incorporate posterior distributions from Bayesian hierarchical models (Chapter~\ref{chap:bayesian-models}):

\begin{itemize}
    \item QB passing yard projections with uncertainty (v2.5 model)
    \item 80 player projections for 2024--2025 seasons
    \item Mean $\pm$ 90\% credible intervals
\end{itemize}

\subsection{Feature Engineering}

\subsubsection{Team-Level Aggregation}

For each team in each game, we engineer features aggregating player statistics from recent games:

\begin{table}[h]
\centering
\caption{Player-Derived Team Features}
\label{tab:player-features}
\begin{tabular}{lll}
\toprule
\textbf{Category} & \textbf{Feature} & \textbf{Calculation} \\
\midrule
\multicolumn{3}{l}{\textit{Offensive Aggregates}} \\
QB Passing & \texttt{qb\_pass\_yards\_avg} & Mean passing yards (recent games) \\
 & \texttt{qb\_pass\_yards\_std} & Std dev (uncertainty measure) \\
 & \texttt{qb\_pass\_tds\_avg} & Mean passing touchdowns \\
 & \texttt{qb\_completion\_pct} & Completions / attempts \\
\midrule
RB Rushing & \texttt{rb\_rush\_yards\_avg} & Mean rushing yards \\
 & \texttt{rb\_yards\_per\_carry} & Rushing yards / carries \\
 & \texttt{rb\_rush\_tds\_avg} & Mean rushing touchdowns \\
\midrule
WR Receiving & \texttt{wr\_rec\_yards\_avg} & Mean receiving yards \\
 & \texttt{wr\_receptions\_avg} & Mean receptions \\
 & \texttt{wr\_rec\_tds\_avg} & Mean receiving touchdowns \\
\midrule
\multicolumn{3}{l}{\textit{Composite Metrics}} \\
Total Offense & \texttt{total\_offense\_yards} & Sum of QB + RB + WR yards \\
Star Power & \texttt{star\_concentration} & Max player yards / total yards \\
\midrule
\multicolumn{3}{l}{\textit{Differentials (Home - Away)}} \\
 & \texttt{diff\_total\_offense} & Home - Away total offense \\
 & \texttt{diff\_qb\_pass\_yards} & Home - Away QB passing \\
 & \texttt{diff\_rb\_rush\_yards} & Home - Away RB rushing \\
\bottomrule
\end{tabular}
\end{table}

\textbf{Rolling Window}: Recent games from current season (week 1--$t-1$); if unavailable, use prior season.

\textbf{Missing Data}: For teams/players with insufficient history, impute league-average values.

\subsubsection{Feature Dimensionality}

Total features per game:
\begin{itemize}
    \item Home team: 11 features
    \item Away team: 11 features
    \item Differentials: 3 features
    \item \textbf{Total}: 33 features
\end{itemize}

This is substantially fewer than GNN embeddings (140K--340K parameters), reducing overfitting risk.

\subsection{Model Architecture}

\subsubsection{XGBoost Classifier}

We train an XGBoost \citep{chen2016xgboost} binary classifier predicting home team win probability:

\begin{align*}
P(\text{Home Win} \mid X) &= \sigma(F(X)) \\
F(X) &= \sum_{k=1}^{K} \eta \cdot f_k(X)
\end{align*}

where:
\begin{itemize}
    \item $X \in \mathbb{R}^{33}$ = feature vector
    \item $f_k$ = decision tree (boosted sequentially)
    \item $\eta$ = learning rate (0.05)
    \item $K$ = number of estimators (100)
    \item $\sigma$ = logistic function
\end{itemize}

\subsubsection{Hyperparameters (Conservative for Small Data)}

\begin{itemize}
    \item \textbf{Max depth}: 3 (shallow trees to prevent overfitting)
    \item \textbf{Learning rate}: 0.05 (slow learning)
    \item \textbf{Subsample}: 0.8 (bootstrap aggregation)
    \item \textbf{Column subsample}: 0.8 (random feature selection)
    \item \textbf{Min child weight}: 5 (regularization)
    \item \textbf{Gamma}: 1.0 (pruning threshold)
\end{itemize}

These conservative settings prioritize generalization over in-sample fit.

\subsection{Training and Evaluation}

\subsubsection{Temporal Validation Strategy}

\begin{enumerate}
    \item \textbf{Training}: 2019--2023 (1,327 games)
    \item \textbf{Validation}: 2024 (272 games)
    \item \textbf{Test}: 2025 (116 games, first 7 weeks)
\end{enumerate}

\textbf{No data leakage}: Features use only games prior to prediction week.

\subsubsection{Results}

Table~\ref{tab:player-to-game-results} presents performance across train, validation, and test sets.

\begin{table}[h]
\centering
\caption{Player-to-Game Model: Performance Summary}
\label{tab:player-to-game-results}
\begin{tabular}{lccc}
\toprule
\textbf{Metric} & \textbf{Train} & \textbf{Val (2024)} & \textbf{Test (2025)} \\
\midrule
Brier Score & 0.1854 & 0.2304 & 0.2439 \\
Log Loss & --- & 0.6523 & 0.6804 \\
Accuracy & 76.3\% & 61.4\% & 55.2\% \\
\midrule
\textit{Comparison} & & & \\
State-Space Brier & 0.2364 & 0.2458 & 0.2615 \\
\textbf{Improvement} & \textbf{+21.6\%} & \textbf{+6.3\%} & \textbf{+6.7\%} \\
\bottomrule
\end{tabular}
\end{table}

\subsubsection{Key Findings}

\begin{enumerate}
    \item \textbf{Validation improvement}: Brier 0.2304 vs.\ state-space 0.2458 ($+6.3\%$)
    \item \textbf{Test improvement}: Brier 0.2439 vs.\ state-space 0.2615 ($+6.7\%$)
    \item \textbf{Consistent generalization}: Unlike GNN/ensemble overfitting, player-to-game shows stable validation$\to$test performance
    \item \textbf{Train-test gap}: 0.1854 $\to$ 0.2439 (31\% degradation), acceptable for sports prediction
\end{enumerate}

\subsection{Feature Importance Analysis}

Table~\ref{tab:feature-importance} shows the top 5 most important features learned by the model.

\begin{table}[h]
\centering
\caption{Top 5 Feature Importance Scores}
\label{tab:feature-importance}
\begin{tabular}{llr}
\toprule
\textbf{Rank} & \textbf{Feature} & \textbf{Importance} \\
\midrule
1 & \texttt{away\_qb\_pass\_tds\_avg} & 0.0591 \\
2 & \texttt{away\_qb\_completion\_pct} & 0.0473 \\
3 & \texttt{away\_wr\_rec\_tds\_avg} & 0.0443 \\
4 & \texttt{home\_qb\_completion\_pct} & 0.0437 \\
5 & \texttt{away\_rb\_rush\_tds\_avg} & 0.0421 \\
\bottomrule
\end{tabular}
\end{table}

\textbf{Insights}:
\begin{itemize}
    \item \textbf{QB dominance}: Passing TDs and completion \% are most predictive
    \item \textbf{Away team bias}: 4 of top 5 features are away team metrics (model learns home-field advantage is already captured elsewhere)
    \item \textbf{Touchdowns matter}: TD rates more important than raw yards
    \item \textbf{Star power not top-5}: Concentration metric ranked lower (rank 15), suggesting depth matters too
\end{itemize}

\subsection{Comparison to Alternative Approaches}

\begin{table}[h]
\centering
\caption{Comparison of Game Prediction Approaches (Test Brier Score)}
\label{tab:approach-comparison}
\begin{tabular}{lccl}
\toprule
\textbf{Model} & \textbf{Val Brier} & \textbf{Test Brier} & \textbf{Status} \\
\midrule
State-Space (Kalman) & 0.2458 & 0.2615 & Robust baseline \\
5-Component Ensemble & 0.1575 & 0.2684 & Severe overfit \\
GNN (Best config) & --- & 0.2500 & 45\% test accuracy \\
\textbf{Player-to-Game} & \textbf{0.2304} & \textbf{0.2439} & \textbf{Best generalization} \\
\bottomrule
\end{tabular}
\end{table}

\textbf{Why Player-to-Game Succeeds Where GNNs Failed}:

\begin{enumerate}
    \item \textbf{Fewer parameters}: 33 features vs.\ 140K--340K GNN parameters
    \item \textbf{Explicit aggregation}: Handcrafted features vs.\ learned message passing
    \item \textbf{Interpretability}: Can inspect which player stats matter
    \item \textbf{Regularization}: XGBoost's conservative hyperparameters prevent overfitting
\end{enumerate}

\subsection{Limitations and Future Work}

\subsubsection{Current Limitations}

\begin{enumerate}
    \item \textbf{No defensive features}: Placeholder values used (league averages); real defensive stats (sack rate, yards allowed vs.\ position) would improve predictions
    \item \textbf{Missing 2025 player data}: Test set uses 2024 player stats as proxy
    \item \textbf{No uncertainty propagation}: Bayesian player projections loaded but not yet integrated into features
    \item \textbf{Static aggregation}: Simple mean/std ignores momentum, injury status, weather
\end{enumerate}

\subsubsection{Promising Extensions}

\textbf{Option 2 from Research Roadmap}: Bayesian Hierarchical Pooling

Rather than XGBoost, use a Bayesian model to properly propagate uncertainty:

\begin{align*}
\text{Team Score} &\sim \mathcal{N}(\mu_{\text{team}}, \sigma^2_{\text{team}}) \\
\mu_{\text{team}} &= f(\theta_{\text{QB}}, \theta_{\text{RB1}}, \theta_{\text{RB2}}, \ldots, \epsilon_{\text{synergy}}) \\
\theta_{\text{player}} &\sim \mathcal{N}(\mu_{\text{player}}, \sigma^2_{\text{player}}) \quad \text{(from BNN posteriors)}
\end{align*}

This would:
\begin{itemize}
    \item Properly account for player projection uncertainty
    \item Learn position-specific weights (e.g., QB = 0.4, RB = 0.2)
    \item Capture team chemistry via $\epsilon_{\text{synergy}}$
\end{itemize}

\textbf{Option 5 from Research Roadmap}: Reinforcement Learning

Use player features as state representations for sequential decision-making:

\begin{itemize}
    \item \textbf{State}: QB completion \%, RB yards/carry, WR separation, fatigue
    \item \textbf{Action}: Play call (run vs.\ pass, formation, target)
    \item \textbf{Reward}: Change in win probability (EPA)
\end{itemize}

Start with tractable problems (4th down decisions) before full game strategy optimization.

\subsection{Implications for Sports Analytics}

\subsubsection{Practical Takeaways}

\begin{enumerate}
    \item \textbf{Simple aggregation beats complex graphs}: Feature engineering (6.7\% improvement) outperforms GNNs (45\% accuracy) despite lower sophistication
    \item \textbf{Player data adds value}: Even basic stats (QB completion \%, RB yards/carry) improve over team ratings alone
    \item \textbf{Regularization is critical}: Conservative XGBoost hyperparameters prevent overfitting on limited data (1,327 games)
    \item \textbf{Interpretability aids trust}: Transparent feature importance (Table~\ref{tab:feature-importance}) enables validation by domain experts
\end{enumerate}

\subsubsection{Betting Market Implications}

Assuming efficient markets (bookmakers use similar features), the 6.7\% Brier improvement translates to:

\begin{itemize}
    \item \textbf{Edge detection}: Better calibration $\Rightarrow$ more accurate identification of mispriced lines
    \item \textbf{Kelly criterion sizing}: Improved win probability estimates $\Rightarrow$ optimal bet sizing
    \item \textbf{Line shopping advantage}: Can compare predictions across multiple books to find best price
\end{itemize}

However, backtests in Section~\ref{sec:game-ensemble} showed that even 76\% historical accuracy does not guarantee profitable deployment (model overfitting, market efficiency, vig).

\subsection{Conclusion}

We demonstrate that aggregating player-level statistics into team-level features improves game prediction beyond team ratings alone. The player-to-game model achieves:

\begin{itemize}
    \item \textbf{6.3\% validation improvement} over state-space baseline (Brier 0.2304 vs.\ 0.2458)
    \item \textbf{6.7\% test improvement} (Brier 0.2439 vs.\ 0.2615)
    \item \textbf{Robust generalization}: Unlike GNNs and ensembles, minimal overfitting
\end{itemize}

This validates the hypothesis that player heterogeneity matters for game outcomes, while affirming that simplicity beats complexity for limited data regimes. The feature engineering approach (Option 1) provides a practical baseline; Bayesian hierarchical pooling (Option 2) and reinforcement learning (Option 5) offer principled next steps for incorporating uncertainty and sequential decision-making.
